\scanpage{191}
% archive sheet number: 190

universel pour les \inserted{$\operatorname{Pro}(\underline C)$-}morphismes d'objets de $\underline C$ en $\unsure{\mathscr X|\mathcal X}$; si $\varprojlim_{\underline C}\mathscr X$ existe pour tt $\mathscr X$, cela signifie précisément que le foncteur \deleted{\unsure{d'inclusion}} pleinement fidèle $\underline C\longrightarrow\operatorname{Pro}(\underline C)$ admet un \deleted{\unsure{droit}}\inserted{adjoint} de Kan, ``projection'' de $\operatorname{Pro}(\underline C)$ \deleted{\unsure{sur}} sur $\underline C$.

On ne confondra pas cette situation avec la suivante: \deleted{\unsure{alors}} il n'y a pas néc. de limites projectives ds $\underline C$, mais on a plongé $\underline C$ comme \uline{sous-catégorie pleine} ds $\underline C'$ telle que \deleted{\unsure{tt|tout}} système projectif de $\underline C$ ait une limite projective \uline{ds $\underline C'$}; alors on a un foncteur naturel
\[
\operatorname{Pro}(\underline C)\longrightarrow\underline C'
\]
qui \uline{parfois} \deleted{\unsure{défini}} (en un cas fréquent) sera pleinement fidèle. [Il en est p.ex. ainsi \unsure{chaque} fois que les objets de $\underline C$ sont artiniens ??]

\uline{Exemple} les pro-(schémas affines de type fini) sur $S$ \deleted{\unsure{soit}} (schéma quasi-compact ou \unsure{préschéma} loc.\unsure{noeth\textsuperscript{n}}) s'identifient exactement aux \uline{schémas} \uline{affines} sur $S$.
% slanted
\margin{N.B. $\operatorname{Pro}(\underline C\text{-grps})\longrightarrow(\operatorname{Pro}\underline C)\text{-grps}$ est un foncteur \uline{pleinement} \uline{fidèle}}

\uline{Question} \deleted{Supposons $S$ \unsure{noethérien}.} Un schéma en groupes affine sur $S$ est-il un pro (groupe affine de type fini), i.e. si $S=\operatorname{Spec}(k)$
